% Post-layout static results. Macros, tables and figure from make_sim_report.py --pex
\input{results_pex.tex}
\section{Post-layout simulation}
\label{sec:postlayout}
The joined netlists of Section~\ref{sec:split} (transistors and coupling capacitances, no resistance)
replace the schematic netlist in the same static characterization: same conditions as in the
summary, staircase and bisection of all 31 transition levels to
\SI{0.39}{mV}. With 68\,519 capacitors on top of the transistors one conversion takes several minutes
instead of a few seconds, so only the static transfer is simulated (taped-out netlist: \SI{7.2}{h}
with 26 jobs). The
sine records are not: at low frequency the dynamic part adds nothing (Section~\ref{sec:dynamic}), and the
ENOB of a sine quantized with the post-layout transition levels (ENOB$_{st}$: a coherent sine of 27 LSB
peak to peak, in the converter's own end-point LSB) is what a low-frequency record would give.

\begin{table}[H]\centering\small
\input{table_pex.tex}
\caption{Static results of the schematic and of the post-layout netlists: the bit weights of a
least-squares fit to the transition levels, in the end-point LSB, the end-point LSB, the INL and the
ENOB of a sine quantized with the transition levels. The last two rows are the taped-out netlist with
only the capacitors that touch a switched or floating plate of the DAC (VSH, FL of every bit), and with
all the other capacitors.}
\label{tab:pex}
\end{table}

\begin{figure}[H]\centering
\includegraphics{fig/pex_dnl_inl.pdf}
\caption{DNL and INL of the schematic and of the two post-layout netlists, each in its own end-point
LSB.}
\label{fig:pex_dnl_inl}
\end{figure}

\subsection{Result}
\begin{itemize}
  \item \textbf{The layout breaks the binary weights.} The taped-out netlist has bit weight errors of
        \TpBitEA{}/\TpBitEB{}/\TpBitEC{}/\TpBitED{}/\TpBitEE{}\,\% (bit 0 to 4), DNL from
        \num{\TpDnlMin} to \num{\TpDnlMax} LSB and \TpInlRms{} LSB rms of INL. The ENOB$_{st}$ falls from
        \SchStEnob{} bits (schematic) to \TpStEnob{} bits. The hierarchical re-layout gives
        \HiBitEA{}/\HiBitEB{}/\HiBitEC{}/\HiBitED{}/\HiBitEE{}\,\%, with the same signs; its INL is smaller
        (\HiInlRms{} LSB rms, ENOB$_{st}$ \HiStEnob{} bits) because its two halves are exact mirrors.
  \item \textbf{The DAC is attenuated.} The end-point LSB grows from \SI{\Lsb}{mV} to \SI{\TpLsb}{mV}
        (taped-out, gain error \SI{\TpGainErr}{\percent}) and \SI{\HiLsb}{mV} (re-layout): the wiring
        capacitance of the top plates (VOUTH/VOUTL to the supplies and to each other) does not take part in
        the redistribution and divides the DAC step.
  \item \textbf{Pattern.} Bit 0 is the most reduced, bits 1 and 3 are smaller and bits 2 and 4 larger
        than binary. The pattern follows the plate wiring (Table~\ref{tab:par_tapeout}): bits 1 and 3 have
        the extra bridges on Metal5/TopMetal1 and about \SI{10}{\percent} more wiring capacitance per unit
        on both plates; the MSB has its switched plate (VSH) coupled \SI{11}{fF} per unit to the top plate
        of its own half (VOUTH for H, VOUTL for L), against about \SI{5}{fF} to each top plate for the other bits; the floating plates
        (FL) of bits 0 to 3 couple 7 to \SI{10}{fF} per unit to the top plate of the other half, the MSB
        less than \SI{1}{fF}. These couplings are of the order of \SI{10}{\percent} of the unit
        capacitance (\SI{67.8}{fF} with FL floating), so they change the effective charge of each bit by
        several percent, differently for each bit.
  \item \textbf{Decomposition.} Keeping only the capacitors on the plates (VSH, FL) gives
        \PlBitEA{}/\PlBitEB{}/\PlBitEC{}/\PlBitED{}/\PlBitEE{}\,\% (INL \PlInlRms{} LSB rms);
        \ifdefined\NpBitEA
        keeping all the others gives \NpBitEA{}/\NpBitEB{}/\NpBitEC{}/\NpBitED{}/\NpBitEE{}\,\%
        (INL \NpInlRms{} LSB rms).
        \else
        the run with all the other capacitors is not finished.
        \fi
        Each part alone already has the pattern of the full netlist; the effects do not simply add
        (the sum of the two is larger than the errors of the full netlist). The attenuation of the DAC comes only from the
        plates (LSB \SI{\PlLsb}{mV} with only the plate capacitors, \SI{\NpLsb}{mV} without them).
        Without the plate capacitors the errors come from the other parasitics around the DAC: the
        capacitance of the top plates and of the switch control nets, which couple the switching of
        each bit into the comparator inputs.
\end{itemize}

\begin{table}[H]\centering\scriptsize
\input{table_par_tapeout.tex}
\caption{Wiring capacitance of the plates of the taped-out layout per unit capacitor (\si{fF}), split
by the net on the other side (\texttt{sim/pex/dacpar.py}). VSH is the switched plate of a bit (to the
input, VOUTH or VOUTL), FL the floating plate. VOUTH and VOUTL are the two comparator inputs (top
plates); H is the lower half, whose MSB connects to VOUTH, L the upper half (MSB on VOUTL, the other bits
with VOUTH and VOUTL exchanged). ``VSH--FL'' is the wiring coupling between the two plates of the same bit, in parallel with the
unit capacitors.}
\label{tab:par_tapeout}
\end{table}
\IfFileExists{table_par_hier.tex}{%
In the hierarchical re-layout the L half is an exact mirror of H (the same numbers with VOUTH and VOUTL
exchanged, Table~\ref{tab:par_hier}). In the taped-out layout it is not: the floating plates of the L
half couple about \SI{10}{fF} to VOUTH and less than \SI{1}{fF} to VOUTL, the H half 3 and \SI{7}{fF}.
This asymmetry between the halves is why the taped-out INL is larger.
\begin{table}[H]\centering\scriptsize
\input{table_par_hier.tex}
\caption{As Table~\ref{tab:par_tapeout}, hierarchical re-layout.}
\label{tab:par_hier}
\end{table}}{}

\subsection{Comparison with the chip}
The chip measures bit weight errors of \num{-7.8}/\num{-5.7}/\num{+6.7}/\num{-8.4}/\num{+3.5}\,\%
(report of the MSO measurement), an INL of about \num{0.5} LSB rms and an ENOB of \num{3.8} bits at low
frequency. The taped-out post-layout netlist has the same sign on all five bits, errors within about
\SI{1}{\percent} of the measured ones on bits 1 and 4 and within \SI{5}{\percent} on the others; it reaches about two thirds of the measured INL and half of the measured
loss of ENOB (from \SchStEnob{} to \TpStEnob{} bits, against \num{3.8}). The layout parasitics of the
DAC are therefore the main cause of the measured nonlinearity. What remains (in the measurement bit 0 is less
reduced, bit 2 larger and bit 3 smaller) is of the size that random mismatch of the unit capacitors, the
references of the board or the parasitic resistance (not extracted) can give.
